\documentclass[10pt,a4paper]{article}
\usepackage[a4paper,left=1.5cm,right=1.5cm,top=1.15cm,bottom=1.05cm]{geometry}
\usepackage{newtxtext,newtxmath}
\usepackage{microtype}
\usepackage{xcolor}
\usepackage{enumitem}
\usepackage{titlesec}
\usepackage{tikz}
\usetikzlibrary{arrows.meta,positioning}
\usepackage[hidelinks]{hyperref}

\definecolor{accent}{HTML}{8C1515}   % Stanford cardinal
\definecolor{tum}{HTML}{0065BD}      % TUM blue
\definecolor{muted}{HTML}{505050}

\titleformat{\section}[runin]{\normalfont\bfseries\color{accent}}{}{0pt}{}[.]
\titlespacing*{\section}{0pt}{2.8pt}{0.5em}
\setlength{\parindent}{0pt}
\setlength{\parskip}{1.3pt}
\setlist{nosep,leftmargin=1.2em,itemsep=0.4pt,topsep=0.6pt}
\linespread{0.97}
\pagestyle{empty}

\newcommand{\hyp}[1]{\textbf{\color{tum}#1}}
\newcommand{\cit}[1]{\textsuperscript{\color{muted}[#1]}}
\newcommand{\step}[1]{\textbf{\color{tum}#1}}

\begin{document}

% ---------------- Header ----------------
{\footnotesize\color{muted} Exposé · CodeX Stanford Scholarship – Master's Thesis Legal Tech (WS 26/27) \hfill Mert Ayvazoglu · TUM School of CIT}\\[-2pt]
{\color{accent}\rule{\linewidth}{0.8pt}}\\[0pt]
{\fontsize{13.2}{15}\selectfont\bfseries Right Outcome, Wrong Reasons? Proof-Grounded Process Rewards\\[-1pt] for Faithful Statutory Reasoning in Large Language Models}\\[1.5pt]
{\small\itshape Rewarding the derivation, not just the answer: can executable law teach legal LLMs to reason faithfully?}\\[1pt]
{\small\textbf{Supervision:} Prof.\ Dr.\ Matthias Grabmair (TUM Legal Tech) \hfill
\textbf{CodeX mentors:} Prof.\ Michael Genesereth (primary) · Dr.\ Marzieh Nabi}\\[-3pt]
{\color{accent}\rule{\linewidth}{0.4pt}}

% ---------------- Motivation ----------------
\section{Motivation}
An LLM that reaches the right tax outcome by applying the wrong rule earns full reward under today's dominant training recipe, reinforcement learning with verifiable rewards (RLVR). A human official who did the same would issue a decision with defective reasons: German authorities must state the grounds ``\emph{that moved them} to their decision'' (§\,39(1) VwVfG; cf.\ Art.\,41(2)(c) CFR), and the AI Act gives persons affected by high-risk legal AI, e.g.\ AI assisting courts in applying the law to facts (Annex~III No.\,8(a)), a right to ``clear and meaningful explanations'' of the decision's main elements (Art.\,86). Such explanations are worth something only if they are \emph{faithful}: the stated reasons must be the operative ones. Legal LLMs already show the gap. Models collapse on integrated hierarchical legal reasoning and spend \emph{more} tokens on wrong answers than on right ones\cit{1}; in 15–52\% of cases they report solver-consistent answers without executing the solver\cit{2}; and their accuracy drops sharply on longer statutory derivations\cit{3}. Meanwhile, legal RL models are being trained on rewards that never verify the derivation\cit{4}. \textbf{Key insight:} \emph{executable law} (Rules as Code) knows not only the answer but the \emph{derivation}, which makes it a process verifier. So far it has served only to evaluate LLMs\cit{3,5} or been proposed as a reward without any training experiments\cit{6}. Process rewards from symbolic oracles help RL in games\cit{7}, but they are untested in law and have never been evaluated for faithfulness.

% ---------------- RQ ----------------
\section{Research question}
\emph{Do outcome rewards train legal LLMs to be right for the wrong reasons, and do process rewards computed from executable-law proof trees make their reasoning faithful?}
\begin{itemize}
  \item[\hyp{H1}] \textbf{Shortcut:} outcome-only RL raises in-distribution accuracy but \emph{lowers} faithfulness; the ``thinking longer, not smarter'' signature\cit{1} is induced by training itself.
  \item[\hyp{H2}] \textbf{Repair:} proof-grounded process rewards reach comparable accuracy with significantly higher faithfulness and reason adequacy.
  \item[\hyp{H3}] \textbf{Transfer:} on held-out statute families, process-RL generalises better than outcome-RL and than imitation (SFT on solver traces).
\end{itemize}

% ---------------- Method ----------------
\section{Methodology}
\begin{center}\vspace{-4pt}
\begin{tikzpicture}[
  node distance=3.2mm,
  box/.style={draw=#1,fill=#1!7,rounded corners=2pt,align=center,font=\scriptsize,inner sep=2pt,minimum height=9mm,text width=29mm},
  arr/.style={-{Stealth[length=1.6mm]},semithick,color=muted}]
  \node[box=tum] (law) {\textbf{Executable statutes}\\SARA tax code (Prolog)\\OpenExempt bankruptcy};
  \node[box=tum,right=of law] (gen) {\textbf{Case generator}\\NL facts + distractors\\+ gold \emph{proof tree}};
  \node[box=accent,right=of gen] (pol) {\textbf{Policy LLM} (4–8B, LoRA)\\free CoT + cited steps\\{[§\,rule $\mapsto$ values]}};
  \node[box=accent,right=of pol] (rew) {\textbf{GRPO reward}\\(a) outcome only vs.\\(b) re-executed step credit};
  \node[box=black,right=of rew] (eval) {\textbf{Reward-disjoint eval}\\counterfactual pairs, CoT\\interventions, transfer};
  \draw[arr] (law)--(gen); \draw[arr] (gen)--(pol); \draw[arr] (pol)--(rew); \draw[arr] (rew)--(eval);
  \draw[arr] (rew.north) to[bend right=14] node[above,font=\tiny,color=muted,inner sep=1pt]{policy update} (pol.north);
\end{tikzpicture}\vspace{-5pt}
\end{center}
\step{(1) Environment.} I wrap existing expert encodings, SARA's Prolog encoding of US tax provisions\cit{5,8} and OpenExempt's bankruptcy-exemption statutes\cit{3}, as generators of unlimited fact patterns. Each pattern is rendered to natural language, paraphrased by an LLM and seeded with distractors. A Prolog meta-interpreter (for SARA) and an instrumented solver (for OpenExempt) return the gold proof tree, \emph{including correctly rejected alternatives}.
\step{(2) Training.} An open reasoning model (e.g.\ Qwen3-4B) is trained with LoRA and GRPO\cit{9}. It reasons freely but cites each rule application inline. Every cited step is \emph{verified by re-execution}: its bound values must reproduce the step in the solver, so the tags cannot be hacked by copying them. Reward (b) credits valid steps, including correctly rejected exceptions, and penalises false ones; the credits are summed per rollout and group-normalised. The core contrast is (a) vs.\ (b), with SFT on solver traces as a control and an autoformalise-and-execute agent as an optional ceiling. Budget: about 300 H100-hours (LRZ), 3 seeds.
\step{(3) Reward-disjoint evaluation.} Faithfulness is scored on the free CoT, never on the tags. \emph{Counterfactual minimal pairs} test whether answers track the legally operative facts: flipping a fact that is necessary in every proof must change the outcome, flipping an irrelevant fact must not, and each pair is verified by re-running the solver. \emph{CoT truncation and corruption}\cit{10} tests whether answers causally depend on the stated reasoning; token spend on wrong versus right answers is compared within difficulty-matched items. \emph{Reason adequacy} turns §\,39 VwVfG into a metric: recall and precision of the essential legal and factual grounds (the necessary proof nodes) stated in the explanation, validated on 200 items against ratings by legal experts. \emph{Transfer} uses leave-one-statute-family-out splits; LegalBench rule-application tasks\cit{11} and hierarchical case-based reasoning\cit{1} serve as exploratory tests on law that has not been formalised.

% ---------------- Contribution ----------------
\section{Contribution \& Legal Tech relevance}
(i)~To my knowledge, the first empirical test of whether RLVR trains legal AI into shortcut reasoners. (ii)~An open RL environment and faithfulness suite built from Rules as Code. (iii)~A measurable version of the duty to give reasons, usable for auditing legal AI. The work turns formalised law, e.g.\ from TUM's Digital-Ready Tax Law project, into \emph{training infrastructure} for trustworthy legal AI.
\textbf{Why CodeX.} With \emph{Genesereth}, I will define step and failed-branch credit over logic-programming proof semantics (Prolog/Epilog). With \emph{Nabi} (System-2 legal reasoning; LLMs encoding insurance contracts as logic programs\cit{12}), I will add her insurance encodings as a third held-out domain (a stretch goal). The CodeX/Stanford Law community supplies expert raters; training runs remotely on LRZ.
\textbf{Plan (6 months).} M1 environment and proof extraction (Munich); M2–4 training and ablations (CodeX); M5 transfer and expert validation; M6 writing. Target: ACL/EMNLP 2027 or ICAIL; code and data released openly.
\textbf{Risks.} If process-GRPO is unstable, I fall back to expert iteration on proof-filtered samples. If H1 is refuted, that is itself a publishable reassurance for legal RLVR, and H2/H3 become questions of sample efficiency and transfer.

% ---------------- References ----------------
\vspace{0pt}{\color{accent}\rule{\linewidth}{0.3pt}}\\[-1.5pt]
{\fontsize{6.8}{7.7}\selectfont\color{muted}
[1]~Zhang, Grabmair, Gray, Ashley: Thinking Longer, Not Always Smarter: Evaluating LLM Capabilities in Hierarchical Legal Reasoning, CSLAW~'26.
[2]~Wang et al.: Know Your Limits: On the Faithfulness of LLMs as Solvers and Autoformalizers in Legal Reasoning, arXiv:2606.16118.
[3]~Servantez et al.: OpenExempt: A Diagnostic Benchmark for Legal Reasoning, Findings of ACL 2026.
[4]~Dai et al.: LegalΔ: Enhancing Legal Reasoning in LLMs via RL with CoT-Guided Information Gain, arXiv:2508.12281.
[5]~Holzenberger, Blair-Stanek, Van Durme: A Dataset for Statutory Reasoning in Tax Law (SARA), NLLP 2020.
[6]~Heydari, Leowald: Closing the Loop: Formally Verified Law as a Reward Signal, arXiv:2606.23913.
[7]~Yuan et al.: Verifiable Process Rewards for Agentic Reasoning, arXiv:2605.10325.
[8]~Jurayj, Holzenberger, Van Durme: Language Models and Logic Programs for Trustworthy Tax Reasoning, AAAI 2026.
[9]~Shao et al.: DeepSeekMath (GRPO), arXiv:2402.03300.
[10]~Lanham et al.: Measuring Faithfulness in Chain-of-Thought Reasoning, arXiv:2307.13702.
[11]~Guha et al.: LegalBench, NeurIPS 2023.
[12]~Kant, Kant, Nabi, Carlson, Ma: Equitable Access to Justice: Logical LLMs Show Promise, arXiv:2410.09904.\par}

\end{document}
